\section{Discussion}
\label{sec:discussion}
% DRAFT (Claude, Nov 27): awaits the author's rewrite. The PFR numbers are from
% the PFR extraction (judge labels, Nov 26); once the cross-area appendix (PR 65)
% is in, point to it with \ref{app:projects}.

Lean proofs in Mathlib mostly stay in one area of mathematics. When a proof
does leave its main area, the off-area steps come in runs, and the share of
steps that move is no larger than judge disagreement alone would give. The
lemmas those steps cite, by contrast, mix areas almost at random. We read
this as a property of the library: Mathlib states cross-area facts as lemmas,
so a proof crosses fields by citing a bridge lemma while its goal stays where
it was. For provers, cross-area knowledge then sits in premise retrieval
rather than in the shape of the goal state. A research formalization known
for joining fields points the same way. In the polynomial Freiman--Ruzsa
proof, judge areas switch more often than in Mathlib, 0.242 times per move,
but against a shuffle of 0.352, a ratio of 0.69 close to Mathlib's 0.71: the
proof mixes more areas and orders them in the same runs.

The label-free model fails for a reason that matters beyond Lean. Goal-state
text differs more between proofs than within one, so any geometry on goal
states, clustered and coarse-grained, recovers proofs and a few long ones,
not areas. Dynamics do not rescue it, whether the estimate is reversible or
not. Areas have to be read from what a step is about, which today takes a
judge, and the judge has to be checked.

Proof depth is a crude clock. Two tactics can do very different amounts of
work, and a single \texttt{simp} or \texttt{linarith} call can close a goal
that would take a human several steps. Switch rates per move therefore
reflect Mathlib's proof style as much as the mathematics itself.

\paragraph{Limitations.}
A judge labels all steps of a proof in one call and gives bookkeeping goals
the surrounding area, so some of the run structure could come from the
judge. A control that shows each proof's steps in shuffled order, with their
positions hidden, separates the two \todo{local and API controls}. The gold
set checks whole statements, not single steps inside a library proof, and
the two main judges come from one family; the local judge from another
family agrees with them only weakly. Premise areas measure what a step uses,
not what it is about, and cover a third of steps. A goal state can hold
several open goals, which makes its area ill-defined after a case split
\todo{share of multi-goal steps}. The traced proofs come from one Mathlib
snapshot and are filtered to 8 or more steps, which may over-represent some
areas. The search runs use a quantized prover at a small budget, and the
prover was trained on data that may overlap with the benchmarks, which
affects every condition alike but limits what absolute numbers mean.
